
\section{Control Unit Subsystem (CUS)}%
\label{sec:cus}

Il CUS costituisce il nodo centrale del sistema: riceve le misure dal TMS, gestisce le richieste degli operatori e determina l'apertura da comandare alla valvola.

\subsection{Architettura software}

Il CUS è sviluppato in Java 17 e utilizza Maven per la gestione del progetto. Il backend si basa su Vert.x per i servizi HTTP e i timer, Eclipse Paho per MQTT e jSerialComm per la comunicazione seriale.

La logica di controllo è concentrata in \code{SystemController}, che gestisce lo stato condiviso e coordina \code{AutomaticPolicy} e \code{MeasurementHistory}, responsabili rispettivamente della politica automatica e dello storico delle misure.
I servizi MQTT, HTTP e seriale gestiscono le comunicazioni con gli altri sottosistemi, come illustrato nella Figura~\ref{fig:cus-components}.

\begin{figure}[htbp]
\centering
\begin{tikzpicture}[node distance=1cm and 1.4cm,block/.append style={font=\footnotesize}]
\node[block] (mqtt) {TmsMqttService\\Parser MQTT};
\node[block,right=of mqtt] (http) {HttpService\\API e file DBS};
\node[block,right=of http] (serial) {WcsSerialAgent\\SerialCommChannel};
\node[block,below=1.2cm of http] (core) {SystemController\\Stato sincronizzato};
\node[block,below left=0.9cm and 0.5cm of core] (policy) {AutomaticPolicy\\Soglie e timer};
\node[block,below right=0.9cm and 0.5cm of core] (history) {MeasurementHistory\\Ultimi 20 campioni};
\draw[flow] (mqtt.south) |- (core.west);
\draw[flow,<->] (http) -- (core);
\draw[flow,<->] (serial.south) |- (core.east);
\draw[flow] (core.south) -- ++(0,-0.3) -| (policy.north);
\draw[flow] (core.south) -- ++(0,-0.3) -| (history.north);
\end{tikzpicture}

\caption{Organizzazione software del CUS e relazioni tra
i componenti principali.}%
\label{fig:cus-components}
\end{figure}

La configurazione viene letta dal file \path{cus.properties}, situato nella directory di avvio dell'applicazione.
In assenza del file vengono utilizzati i valori predefiniti.
Il percorso della directory DBS è relativo alla directory di esecuzione.

I principali parametri del CUS sono riportati nella Tabella~\ref{tab:cus-parameters}.
Le soglie, gli intervalli temporali T1 e T2 e la porta HTTP sono configurabili; il tick periodico e la capacità dello storico sono invece definiti nel codice.

\begin{table}[htbp]
\centering
\begin{tabularx}{\linewidth}{@{}l l X@{}}
\toprule
\textbf{Parametro} & \textbf{Valore} & \textbf{Descrizione} \\
\midrule
\(L_1,L_2\) & \(40\%,70\%\) & Soglie del livello dell'acqua. \\
\(T_1\) & \(5\,\mathrm{s}\) & Intervallo di attesa nella fascia intermedia. \\
\(T_2\) & \(10\,\mathrm{s}\) & Tempo massimo senza misure valide. \\
Tick & \(100\,\mathrm{ms}\) & Intervallo di controllo periodico. \\
\(N\) & 20 & Capacità dello storico delle misure. \\
HTTP & 8080 & Porta del servizio HTTP. \\
\bottomrule
\end{tabularx}
\caption{Parametri principali del CUS e relativi valori
predefiniti o definiti nel codice.}%
\label{tab:cus-parameters}
\end{table}
\FloatBarrier%


\subsection{Gestione dello stato e politica automatica}%
\label{subsec:cus-policy}
All'avvio il sistema si trova in modalità \code{AUTOMATIC}, con apertura della valvola allo \(0\%\) e livello dell'acqua non ancora disponibile.
Il conteggio di T2 inizia all'avvio, in attesa della prima misura valida.


\subsubsection{Politica a stati}
La politica automatica è organizzata attraverso una FSM con cinque stati interni: \code{UNKNOWN}, stato  iniziale o successivo a un reset; \code{LOW}, associato al livello basso; \code{MID\_WAITING}, che rappresenta l'attesa T1; \code{MID\_ACTIVE}, con apertura al 50\%; e \code{HIGH}, con apertura al 100\%.

Le condizioni applicate alle misure valide sono riassunte nella Tabella~\ref{tab:automatic-policy}.

\begin{table}[htbp]
\centering
\begin{tabularx}{\linewidth}{@{}l X@{}}
\toprule
\textbf{Livello} & \textbf{Comportamento} \\
\midrule

\(p<L_1\)
& Passa a \code{LOW} e imposta l'apertura allo \(0\%\). \\

\(p=L_1\)
& Mantiene il \(50\%\) se si trova già in \code{MID\_ACTIVE}; altrimenti passa a \code{LOW} e chiude la valvola. \\

\(L_1<p<L_2\)
& Mantiene il \(50\%\) se già in \code{MID\_ACTIVE}; altrimenti entra in \code{MID\_WAITING}, se necessario, conservando l'apertura precedente durante T1. \\

\(p=L_2\)
& Mantiene il \(100\%\) se già in \code{HIGH}; altrimenti segue la gestione della fascia intermedia. \\

\(p>L_2\)
& Passa a \code{HIGH} e imposta immediatamente l'apertura al \(100\%\). \\

\bottomrule
\end{tabularx}
\caption{Comportamento della politica automatica in funzione del livello e dello stato precedente.}%
\label{tab:automatic-policy}
\end{table}

In \code{MID\_WAITING}, al raggiungimento di T1, la politica passa a \code{MID\_ACTIVE} e imposta l'apertura al \(50\%\). La scadenza viene verificata periodicamente dal CUS, anche in assenza di nuove misure.

Durante l'attesa viene mantenuta l'apertura precedente.
Ad esempio, dopo una discesa da \code{HIGH} nella fascia intermedia, il \(100\%\) rimane attivo fino alla scadenza di T1, quando l'apertura passa al \(50\%\).

Le principali transizioni della politica automatica sono rappresentate nella Figura~\ref{fig:cus-policy}.

\begin{figure}[htbp]
\centering
\begin{tikzpicture}[x=1cm,y=1cm]
\node[state] (unknown) at (0,0) {UNKNOWN\\Politica azzerata};
\node[state,text width=3cm] (low) at (-4.3,-2.1) {LOW\\Apertura 0\%};
\node[state,text width=3cm] (high) at (4.3,-2.1) {HIGH\\Apertura 100\%};
\node[state] (wait) at (0,-4.1) {MID\_WAITING\\Mantiene apertura};
\node[state] (active) at (0,-6.1) {MID\_ACTIVE\\Apertura 50\%};
\draw[flow] (unknown) -- node[right,font=\footnotesize]{\(L_1<p\leq L_2\)} (wait);
\draw[flow] (low.south) |- node[pos=0.3,left,font=\footnotesize]{\(L_1<p\leq L_2\)} (wait.west);
\draw[flow] (high.south) |- node[pos=0.3,right,font=\footnotesize]{\(L_1<p<L_2\)} (wait.east);
\draw[flow] (wait) -- node[right,font=\footnotesize]{\(\Delta t\geq T_1\)} (active);
\end{tikzpicture}
\par\smallskip


\begin{minipage}{0.94\linewidth}
\footnotesize
\raggedright

\textbf{Condizioni e transizioni aggiuntive:}

\begin{itemize}
    \setlength{\itemsep}{3pt}
    \setlength{\parsep}{0pt}
    \setlength{\topsep}{4pt}

    \item \textbf{Soglia L1:}
    con \(p<L_1\) si passa a LOW. A \(p=L_1\), MID\_ACTIVE mantiene il 50\%, mentre gli altri stati passano a LOW.

    \item \textbf{Soglia L2:}
    con \(p>L_2\) si passa a HIGH, che rimane attivo anche a \(p=L_2\).

    \item \textbf{Timer T1:}
    alla scadenza, MID\_WAITING passa a MID\_ACTIVE, anche senza una nuova misura.

    \item \textbf{Reset:}
    riporta la politica allo stato UNKNOWN.
\end{itemize}
\end{minipage}



\caption{Stati e principali transizioni della politica
automatica.}
\label{fig:cus-policy}
\end{figure}


\subsubsection{Misure non valide, timeout e recupero}
Quando il TMS segnala una lettura non valida tramite \code{NO\_ECHO}, il CUS registra un campione \code{null} nello storico e rende il livello corrente non disponibile.
In modalità \code{AUTOMATIC}, la politica viene azzerata, mantenendo invariata l'apertura della valvola.

Il timer T2 viene aggiornato soltanto alla ricezione di misure valide.
Se non ne arrivano per almeno T2, il CUS passa allo stato \code{UNCONNECTED} e imposta l'apertura al \(5\%\).
Questo valore è stato scelto con l'intento di consentire un deflusso limitato ed evitare uno svuotamento rapido della vasca.

Alla prima misura valida dopo la scadenza di T2, il CUS ripristina la modalità \code{AUTOMATIC} e rivaluta il livello.
La rivalutazione parte internamente dall'apertura di fallback al \(5\%\), senza necessariamente trasmettere al WCS un comando intermedio a tale valore.


\subsubsection{Comandi manuali e storico}
Il passaggio alla modalità \code{MANUAL} mantiene l'apertura corrente della valvola.
L'operatore può impostare un'apertura compresa tra \(0\%\) e \(100\%\), tramite l'interfaccia locale o la dashboard.
Le misure del TMS continuano ad aggiornare il livello e lo storico, senza modificare l'apertura comandata.

Il ritorno alla modalità \code{AUTOMATIC} azzera la politica di controllo e rivaluta l'ultima misura, se disponibile.
Durante lo stato \code{UNCONNECTED}, le richieste di cambio modalità e di regolazione dell'apertura vengono rifiutate.

Lo storico conserva in memoria gli ultimi \(N\) campioni ricevuti e accettati, comprese le misure non valide rappresentate da \code{null}.
I campioni sono ordinati cronologicamente in base alla ricezione da parte del CUS.
Lo storico non è persistente e viene azzerato al riavvio.


\subsection{Comunicazione e sincronizzazione}
Il CUS gestisce eventi provenienti da MQTT, HTTP, comunicazione seriale e timer periodici, che possono essere elaborati concorrentemente.
L'accesso allo stato condiviso, alla politica automatica e allo storico è protetto da un lock, mentre le operazioni di comunicazione vengono eseguite al di fuori delle sezioni critiche.

Lo stato del sistema è rappresentato attraverso snapshot immutabili, utilizzati per fornire informazioni coerenti ai componenti esterni.

Gli aggiornamenti destinati al WCS vengono accorpati per privilegiare lo stato più recente, evitando l'accumulo di stati intermedi.
Le richieste ricevute tramite seriale vengono invece elaborate dal controller, che ne verifica la compatibilità con lo stato corrente e restituisce lo stato aggiornato.




\subsection{Validazione, test e limitazioni}%
\label{subsec:cus-validation}
Il CUS è stato verificato attraverso test unitari JUnit, dedicati alla politica automatica, alla gestione dello stato e all'interpretazione dei messaggi MQTT.

I test coprono le transizioni tra modalità operative, il comportamento alle soglie L1 e L2, le scadenze T1 e T2, la gestione delle misure non valide e dei comandi manuali.
Sono inoltre verificate la sincronizzazione degli aggiornamenti destinati al WCS e la coerenza del recupero dopo T2.
Tutti i test unitari eseguiti sono stati superati.

Le verifiche unitarie non comprendono il funzionamento integrato con broker MQTT, dispositivi fisici e dashboard,che richiede prove separate.

Tra le limitazioni del prototipo, il CUS non dispone di un riscontro sulla posizione effettiva della valvola.
La comunicazione seriale non prevede meccanismi automatici di riconnessione e ritrasmissione dei messaggi non
consegnati.
I servizi HTTP e MQTT non implementano autenticazione o cifratura delle comunicazioni.
